\textbf{(i)} $\Rightarrow$ \textbf{(ii)}. If $\mathfrak p\subseteq\mathfrak m$ are prime ideals and $\mathfrak m$ is maximal, both give prime ideals of $(A/\mathfrak N)_{\mathfrak m}$. This ring is a field by \exref{3.10}, so the prime correspondence gives $\mathfrak p=\mathfrak m$.

\textbf{(ii)} $\Leftrightarrow$ \textbf{(iii)} follows from $\overline{\{\mathfrak p\}}=V(\mathfrak p)$.

\textbf{(ii)} $\Rightarrow$ \textbf{(iv)}. Given distinct primes $\mathfrak p,\mathfrak q$, choose $f\in\mathfrak p\setminus\mathfrak q$. The local ring $A_{\mathfrak p}$ has only one prime ideal, so $f/1$ is nilpotent. Thus $sf^n=0$ for some $s\notin\mathfrak p$ and $n\ge1$. The neighborhoods $D(s)$ of $\mathfrak p$ and $D(f)$ of $\mathfrak q$ are disjoint, since $sf$ is nilpotent.

\textbf{(iv)} $\Rightarrow$ \textbf{(iii)} is immediate. To obtain \textbf{(i)} from \textbf{(ii)}, each $(A/\mathfrak N)_{\mathfrak m}$ is reduced and has only one prime ideal. Its maximal ideal is therefore its zero nilradical, so it is a field. Apply \exref{3.10}.

The spectrum is quasi-compact, hence compact when Hausdorff. Each $D(f)$ is quasi-compact and therefore closed in this Hausdorff space. Any two distinct points can be separated by such a clopen set, so no connected subset contains two distinct points.
